\chapter{Problem Definition and System Overview}\label{ch:problem}

\section{Formal Problem Definition}

Let $\mathcal{K}$ be a set of $K=23$ tactical categories (fork, pin, skewer, mating pattern, rook
endgame, and so on; the full list is in Appendix~\ref{app:constants}). A \emph{puzzle} $p$ has a
category $k(p)\in\mathcal{K}$, a starting position, a unique solution line, and a difficulty
$b_p\in\mathbb{R}$ on a logit scale. A \emph{player} $u$ has an overall tactical ability $\theta_u$
and a category-specific offset $\delta_{u,k}$ for each $k\in\mathcal{K}$. The probability that $u$
solves $p$ is modelled as

\begin{equation}\label{eq:rasch}
  P(y_{u,p}=1) \;=\; \sigma\!\left(\theta_u + \delta_{u,k(p)} - b_p\right),
  \qquad \sigma(z) = \frac{1}{1+e^{-z}} .
\end{equation}

\noindent This is the Rasch form used by Elo and Glicko. With ratings, $b_p=(r_p-1500)/173.7178$, so
a 400-point gap corresponds to 10:1 odds.

\paragraph{The decision problem.} At each step $t$ the system chooses a category $k_t$ and a
difficulty window, serves a puzzle $p_t$, and observes $y_t\in\{0,1\}$. The \emph{weakness-targeting}
objective is to concentrate practice on the categories with the lowest $\delta_{u,k}$, measured by
the per-step regret

\begin{equation}\label{eq:regret}
  \mathrm{regret}_t \;=\; \delta_{u,k_t} - \min_{k}\delta_{u,k} \;\geq\; 0 .
\end{equation}

\noindent The \emph{difficulty} objective keeps $P(y_t=1)$ in a productive range. It should be
neither trivially easy nor frustrating; this dissertation targets 0.65 and treats the band
$[0.5,0.85]$ as productive. The two objectives interact. The $\delta$ values are learned from
outcomes, and outcomes depend on the difficulty chosen.

\paragraph{Evidence before the first puzzle.} Before any puzzle is served, the player's recent games
provide \emph{opportunities}: critical positions where one move is clearly best. Each is labelled with
a category and recorded as found or missed. They inform a prior over $(\theta_u,\delta_{u,\cdot})$.
How much they should be trusted is an empirical question, answered in
Section~\ref{sec:cohort-results}.

\section{Conceptual Model: Player--Puzzle--Learning Loop}

Figure~\ref{fig:loop} shows the loop. Diagnosis from games supplies a prior. Selection chooses what
to practise and how hard it should be. Each attempt updates the belief about the player. The player
also changes, since practice produces learning, so the belief must be able to track a moving
target.

\begin{figure}[H]
\centering
\begin{tikzpicture}[
  box/.style={draw, rounded corners, align=center, minimum width=3.2cm, minimum height=1.1cm, font=\small},
  arr/.style={-{Stealth[length=2.5mm]}, thick}]
  \node[box, fill=gray!10] (games) {Player's recent\\games};
  \node[box, fill=blue!8, right=1.2cm of games] (diag) {Diagnosis\\(prior on $\theta,\delta$)};
  \node[box, fill=blue!8, right=1.2cm of diag] (select) {Selection\\(category $+$ difficulty)};
  \node[box, fill=orange!10, below=1.4cm of select] (attempt) {Puzzle attempt\\(solve / fail)};
  \node[box, fill=blue!8, below=1.4cm of diag] (update) {Belief update\\(posterior on $\theta,\delta$)};
  \node[box, fill=green!8, below=1.4cm of games] (learn) {Player learns\\($\delta$ changes)};
  \draw[arr] (games) -- (diag);
  \draw[arr] (diag) -- (select);
  \draw[arr] (select) -- (attempt);
  \draw[arr] (attempt) -- (update);
  \draw[arr] (update) -- (select);
  \draw[arr, dashed] (attempt) to[bend left=20] (learn);
  \draw[arr, dashed] (learn) -- (games);
\end{tikzpicture}
\caption{The player--puzzle--learning loop. Solid arrows are the system's own inference; dashed arrows
are the change in the player that the system must track.}\label{fig:loop}
\end{figure}

\section{Requirements and Design Goals}

\paragraph{Functional requirements.}
\begin{enumerate}[label=FR\arabic*]
  \item Fetch a player's public games and profile from Chess.com without authentication.
  \item Analyse every move with Stockfish; grade errors and record critical positions (opportunities)
        with their tactical category.
  \item Build a weakness profile and priors for the recommender.
  \item Mine puzzles with a unique solution from the player's own games.
  \item Serve puzzles adaptively, choosing category and difficulty, and update after every attempt.
  \item Log each attempt with the serving policy, the selection probability and the pre-outcome
        prediction, so that the system can be evaluated later.
\end{enumerate}

\paragraph{Non-functional requirements and design goals.}
\begin{itemize}
  \item \textbf{Reproducibility.} Engine analysis is limited by node count, not time, so the same game
        always gives the same result. Simulations use fixed seeds and common random numbers.
  \item \textbf{Testability.} The whole test suite runs offline, with no engine or network (231 tests).
  \item \textbf{Graceful degradation.} If Stockfish is missing, or a category has no puzzle in range,
        the system falls back instead of failing.
  \item \textbf{Privacy.} Research data is pseudonymised and only aggregates are published.
  \item \textbf{Honest evaluation.} Every numeric claim can be regenerated from a result file by a
        script.
\end{itemize}

\section{High-Level System Architecture}

Figure~\ref{fig:arch} shows the components and data flow.

\begin{figure}[H]
\centering
\begin{tikzpicture}[
  comp/.style={draw, rounded corners, align=center, minimum width=3.2cm, minimum height=1.0cm, font=\footnotesize},
  arr/.style={-{Stealth[length=2mm]}, thick},
  lab/.style={font=\scriptsize, fill=white, inner sep=1pt}]
  % three columns (x = 0, 4.2, 8.4) and three rows (y = 0, -2.2, -4.4)
  \node[comp, fill=gray!10]   (cc)    at (0,0)      {Chess.com\\Published-Data API};
  \node[comp, fill=blue!8]    (an)    at (4.2,0)    {Game analyser\\(Stockfish + tactic tagger)};
  \node[comp, fill=blue!8]    (prof)  at (8.4,0)    {Profiler\\(opportunities, norms)};
  \node[comp, fill=gray!10]   (pool)  at (0,-2.2)   {Lichess pool\\(5.9M puzzles)};
  \node[comp, fill=blue!8]    (mine)  at (4.2,-2.2) {Puzzle miner\\(4 engine gates)};
  \node[comp, fill=orange!12] (irt)   at (8.4,-2.2) {IRT learner\\(Thompson Sampling)};
  \node[comp, fill=blue!8]    (serve) at (4.2,-4.4) {Serving cascade\\(Flask API)};
  \node[comp, fill=green!8]   (ui)    at (8.4,-4.4) {Player\\(browser SPA)};
  \draw[arr] (cc) -- node[lab, above]{games} (an);
  \draw[arr] (an) -- (prof);
  \draw[arr] (cc) -- node[lab, sloped, below]{own games} (mine);
  \draw[arr] (prof) -- node[lab, right]{priors} (irt);
  \draw[arr] (mine) -- (serve);
  \draw[arr] (pool) -- (serve);
  \draw[arr] (irt) -- node[lab, sloped, above]{category, difficulty} (serve);
  \draw[arr] (serve) -- node[lab, above]{puzzle} (ui);
  \draw[arr] (ui) -- node[lab, right]{solve / fail} (irt);
\end{tikzpicture}
\caption{System architecture. The IRT learner decides the category and difficulty of the next puzzle,
and the serving cascade finds an unseen matching puzzle. Every attempt updates the learner and is
logged with the serving policy, the selection probability and the pre-outcome
prediction.}\label{fig:arch}
\end{figure}

The \textbf{game analyser} evaluates every player move after the opening with Stockfish. It grades
the move by the loss of win probability it causes, and records \emph{opportunities}, meaning
positions where the best move beats the runner-up by at least 10 percentage points of win
probability. The \textbf{tactic tagger} labels the engine's best line with one of the 23 categories.
The \textbf{profiler} turns found and missed opportunities into per-category estimates, shrunk
towards population norms. The \textbf{IRT learner} maintains a Gaussian posterior over $\theta_u$ and
the 23 offsets $\delta_{u,k}$. It chooses a category by Thompson Sampling, and a target difficulty
that gives a predicted 65\% solve rate. The \textbf{serving cascade} finds an unseen puzzle in that
category and rating window. It prefers puzzles mined from the player's own games, and widens its
constraints one at a time when nothing matches. The \textbf{miner} extracts puzzles from the player's
own losses through four verification gates.

\section{Chapter Summary}

The problem is formalised as sequential decision-making under uncertainty, based on a Rasch-type
model with a category-specific ability offset. The system implements that model online, with
evidence from games as a prior and puzzle attempts as the main source of information. The
requirements stress reproducibility, offline testability and evaluation that can be checked, and
these shape the methodology in the next chapter.
